%----------------------------------------------------------------------
\section{Meta-analysis: the technique matrix}\label{app:meta}
%----------------------------------------------------------------------

This appendix records the structure discovered by
profiling the paper's 15 non-trivial formal objects
against six closure techniques specific to
scheduler verification. Following the convention
of [DAF] Appendix D, the analysis is conducted
\emph{after} the spine is complete; it reports
what proof technology is actually load-bearing.

\subsection{The six techniques}

Every closure in the paper uses one or more of:

\begin{description}[nosep,leftmargin=1.5cm]
  \item[S1: $\GF(2)$ linearity] The
    characteristic-2 algebra from [DAF]: $x \oplus
    x = 0$, $-x = x$, linearity of $\oplus$ through
    any linear operator. Used by the XOR-bridge and
    the encoding commutation.
  \item[S2: Block-diagonal decomposition] The matrix
    identity that a block-diagonal matrix-vector
    product computes each block independently. Used
    by \texttt{fused\_matrix\_apply} and by the
    mask-disjointness argument.
  \item[S3: XOR reversibility] The double-XOR
    identity $x \oplus m \oplus m = x$. Used by the
    XOR-bridge's save-restore mechanism.
  \item[S4: Pareto dominance] The partial-order
    fact that if $c \leq c'$ componentwise and
    $c \neq c'$, then $c'$ can be discarded. Used
    by chart-item pruning, Pareto frontier
    extraction, and portfolio construction.
  \item[S5: A$^*$ admissibility] The AI-search fact
    that admissible monotone heuristics make A$^*$
    optimal and complete. Used by the cost-bounded
    Earley completer.
  \item[S6: LP relaxation \& duality] Strong duality
    for linear programs, the envelope theorem for
    sensitivity, and the integrality gap bound.
    Used by the Lagrangian dual and the dual-price
    diagnostic.
\end{description}

\subsection{The closure--technique matrix}

\begin{center}
\scriptsize
\begin{tabular}{l|cccccc|c}
 & S1 & S2 & S3 & S4 & S5 & S6 & Cat \\
\hline
Zeta involutivity (\ref{prop:zeta-involution}) &
  $\bullet$ & & & & & & I \\
Encoding commutation (\ref{prop:encoding-commute}) &
  $\bullet$ & & & & & & I \\
LCM density (\ref{prop:lcm-density}) &
  & & & & & & C \\
XOR-bridge correctness (\ref{prop:xor-bridge}) &
  $\bullet$ & & $\bullet$ & & & & V \\
Fused-apply correctness (\ref{prop:fused-correctness}) &
  $\bullet$ & $\bullet$ & & & & & V \\
Mask disjointness (\ref{prop:mask-disjoint}) &
  & $\bullet$ & & & & & A \\
Time-share commutativity (\ref{prop:ts-commute}) &
  $\bullet$ & & $\bullet$ & & & & A \\
Ordered time-share (\ref{prop:ts-ordered}) &
  $\bullet$ & & $\bullet$ & & & & A \\
Grammar soundness (\ref{thm:soundness}) &
  $\bullet$ & $\bullet$ & $\bullet$ & & & & T \\
Cost-bounded Earley (\ref{thm:earley-complete}) &
  & & & $\bullet$ & $\bullet$ & & T \\
Heuristic monotone (\ref{prop:heuristic-monotone}) &
  & & & & $\bullet$ & & A \\
Complexity (\ref{thm:complexity}) &
  & & & $\bullet$ & & & T \\
Incremental replan (\ref{prop:incremental}) &
  & & & & $\bullet$ & & A \\
Lagrangian (\ref{thm:lagrangian}) &
  & & & & & $\bullet$ & T \\
Pareto extraction (\ref{prop:pareto}) &
  & & & $\bullet$ & & & A \\
Parse-packing (\ref{thm:parse-packing}) &
  $\bullet$ & $\bullet$ & & & & & T \\
Plan transitivity (\ref{prop:plan-trans}) &
  & & & & & & A \\
Register bound (\ref{thm:reg-bound}) &
  & & & & & & T \\
\end{tabular}
\end{center}

Legend: I = identity, V = verification,
A = auxiliary proposition, T = theorem, C = counting.
Technique counts across 18 formal objects:
S1 (GF(2) linearity) \emph{9 uses},
S2 (block-diagonal) \emph{4},
S3 (XOR reversibility) \emph{3},
S4 (Pareto) \emph{3},
S5 (A$^*$) \emph{3},
S6 (LP duality) \emph{1}.

\subsection{Equivalence classes}

\emph{$\{S1\}$-only class}: zeta involution, encoding
commutation. Pure $\GF(2)$ linearity; no
scheduling-level structure. These are the results
inherited most directly from [DAF]'s algebra.

\emph{$\{S1, S3\}$ class}: XOR-bridge correctness,
time-share commutativity, ordered time-share. The
specifically \emph{reversible} uses of $\GF(2)$
linearity.

\emph{$\{S1, S2\}$ class}: fused-apply correctness,
parse-packing. The specifically \emph{parallel}
uses of $\GF(2)$ linearity.

\emph{$\{S1, S2, S3\}$ class}: grammar soundness
alone. The soundness theorem draws on all three
\emph{linearity-adjacent} techniques because it
must cover both spatial and temporal production
families in one induction.

\emph{$\{S4, S5\}$ class}: cost-bounded Earley
completeness. The parser correctness argument is
pure AI-search theory; no CSTZ-specific algebra is
invoked.

\emph{$\{S6\}$ class}: Lagrangian achievability
alone. The LP-duality machinery sits entirely
outside the $\GF(2)$ story.

\subsection{Interpretation}

Two observations fall out of the matrix:

\begin{itemize}[nosep]
  \item \textbf{S1 ($\GF(2)$ linearity) is the
    paper's spine}: it appears in every
    correctness claim about compound productions.
    This reflects the design choice of building
    every compound production on linear atoms; the
    moment the grammar admits a non-linear
    production (beyond the isolated
    $\mathrm{pcnt}_2$), the spine would fracture.
  \item \textbf{S6 (LP duality) is a leaf}: it is
    invoked only by the Lagrangian theorem, which
    stands alone in the technique graph. This
    suggests the Lagrangian dual is a modular
    addition: one could remove it and retain the
    rest of the planner (with a fixed scalar
    objective instead of a budget-constrained
    formulation). Whether the planner would
    then be operationally useful is a separate
    question.
\end{itemize}

The matrix also exhibits a clean factorization
into \emph{algebra} (S1, S2, S3) and
\emph{combinatorial optimization} (S4, S5, S6),
with no claim drawing from both. That is precisely
the architectural separation between the grammar
and the parser: the grammar is the algebra, the
parser is the search. The paper's structural
coherence derives from keeping these two
halves independent.
